% Options for packages loaded elsewhere
\PassOptionsToPackage{unicode}{hyperref}
\PassOptionsToPackage{hyphens}{url}
\documentclass[
  ignorenonframetext,
]{beamer}
\newif\ifbibliography
\usepackage{pgfpages}
\setbeamertemplate{caption}[numbered]
\setbeamertemplate{caption label separator}{: }
\setbeamercolor{caption name}{fg=normal text.fg}
\beamertemplatenavigationsymbolsempty
% remove section numbering
\setbeamertemplate{part page}{
  \centering
  \begin{beamercolorbox}[sep=16pt,center]{part title}
    \usebeamerfont{part title}\insertpart\par
  \end{beamercolorbox}
}
\setbeamertemplate{section page}{
  \centering
  \begin{beamercolorbox}[sep=12pt,center]{section title}
    \usebeamerfont{section title}\insertsection\par
  \end{beamercolorbox}
}
\setbeamertemplate{subsection page}{
  \centering
  \begin{beamercolorbox}[sep=8pt,center]{subsection title}
    \usebeamerfont{subsection title}\insertsubsection\par
  \end{beamercolorbox}
}
% Prevent slide breaks in the middle of a paragraph
\widowpenalties 1 10000
\raggedbottom
\AtBeginPart{
  \frame{\partpage}
}
\AtBeginSection{
  \ifbibliography
  \else
    \frame{\sectionpage}
  \fi
}
\AtBeginSubsection{
  \frame{\subsectionpage}
}
\usepackage{iftex}
\ifPDFTeX
  \usepackage[T1]{fontenc}
  \usepackage[utf8]{inputenc}
  \usepackage{textcomp} % provide euro and other symbols
\else % if luatex or xetex
  \usepackage{unicode-math} % this also loads fontspec
  \defaultfontfeatures{Scale=MatchLowercase}
  \defaultfontfeatures[\rmfamily]{Ligatures=TeX,Scale=1}
\fi
\usepackage{lmodern}
\ifPDFTeX\else
  % xetex/luatex font selection
\fi
% Use upquote if available, for straight quotes in verbatim environments
\IfFileExists{upquote.sty}{\usepackage{upquote}}{}
\IfFileExists{microtype.sty}{% use microtype if available
  \usepackage[]{microtype}
  \UseMicrotypeSet[protrusion]{basicmath} % disable protrusion for tt fonts
}{}
\makeatletter
\@ifundefined{KOMAClassName}{% if non-KOMA class
  \IfFileExists{parskip.sty}{%
    \usepackage{parskip}
  }{% else
    \setlength{\parindent}{0pt}
    \setlength{\parskip}{6pt plus 2pt minus 1pt}}
}{% if KOMA class
  \KOMAoptions{parskip=half}}
\makeatother
\usepackage{longtable,booktabs,array}
\newcounter{none} % for unnumbered tables
\usepackage{calc} % for calculating minipage widths
\usepackage{caption}
% Make caption package work with longtable
\makeatletter
\def\fnum@table{\tablename~\thetable}
\makeatother
\setlength{\emergencystretch}{3em} % prevent overfull lines
\providecommand{\tightlist}{%
  \setlength{\itemsep}{0pt}\setlength{\parskip}{0pt}}
\usepackage{bookmark}
\IfFileExists{xurl.sty}{\usepackage{xurl}}{} % add URL line breaks if available
\urlstyle{same}
\hypersetup{
  hidelinks,
  pdfcreator={LaTeX via pandoc}}

\author{\texorpdfstring{}{}}
\date{}

\begin{document}

\section{Appendix: AI Protocols, Diagrammatic Reasoning, and Notation
Conventions}\label{appendix-ai-protocols-diagrammatic-reasoning-and-notation-conventions}

\begin{frame}{H. AI and Communication Protocols}
\protect\phantomsection\label{h.-ai-and-communication-protocols}
{\def\LTcaptype{none} % do not increment counter
\begin{longtable}[]{@{}
  >{\raggedright\arraybackslash}p{(\linewidth - 6\tabcolsep) * \real{0.2222}}
  >{\centering\arraybackslash}p{(\linewidth - 6\tabcolsep) * \real{0.2778}}
  >{\raggedright\arraybackslash}p{(\linewidth - 6\tabcolsep) * \real{0.2222}}
  >{\centering\arraybackslash}p{(\linewidth - 6\tabcolsep) * \real{0.2778}}@{}}
\toprule\noalign{}
\begin{minipage}[b]{\linewidth}\raggedright
Term
\end{minipage} & \begin{minipage}[b]{\linewidth}\centering
Symbol
\end{minipage} & \begin{minipage}[b]{\linewidth}\raggedright
Definition
\end{minipage} & \begin{minipage}[b]{\linewidth}\centering
First Use
\end{minipage} \\
\midrule\noalign{}
\endhead
A2A Protocol & --- & Google's Agent-to-Agent protocol for
cross-framework agent communication via HTTP/JSON-RPC & §9 \\
ACP & --- & Agent Communication Protocol; standardizes messaging formats
across agents, apps, and humans & §9 \\
ANP & --- & Agent Network Protocol; three-layer architecture for trusted
distributed agent interaction & §9 \\
Categorical deep learning & --- & Deep learning approached through the
lens of category theory (Gavranović et al.) & §9 \\
Double Categorical Systems Theory (DCST) & --- & Framework using
2-categories (horizontal + vertical composition) for explainable
autonomous AI & §9 \\
Functorial encryption & --- & Semantic cryptography: applying a secret
functor to map plaintext categories into ciphertext categories & §9 \\
lambeq & --- & Quantum Natural Language Processing pipeline compiling
DisCoCat diagrams to quantum circuits & §9 \\
MCP & --- & Model Context Protocol; standardizes how AI agents access
external tools and data sources & §9 \\
Parameterized optics / lenses & --- & Categorical constructions modeling
neural network components (Gavranović); attention heads as optics &
§9 \\
QNLP & --- & Quantum Natural Language Processing; quantum computation on
DisCoCat sentence diagrams & §9 \\
Semantic cryptography & --- & Encrypting compositional meaning
structures (functorial encryption, diagram obfuscation, weight masking)
& §9 \\
\bottomrule\noalign{}
\end{longtable}
}
\end{frame}

\begin{frame}{I. Diagrammatic Reasoning}
\protect\phantomsection\label{i.-diagrammatic-reasoning}
{\def\LTcaptype{none} % do not increment counter
\begin{longtable}[]{@{}
  >{\raggedright\arraybackslash}p{(\linewidth - 6\tabcolsep) * \real{0.2222}}
  >{\centering\arraybackslash}p{(\linewidth - 6\tabcolsep) * \real{0.2778}}
  >{\raggedright\arraybackslash}p{(\linewidth - 6\tabcolsep) * \real{0.2222}}
  >{\centering\arraybackslash}p{(\linewidth - 6\tabcolsep) * \real{0.2778}}@{}}
\toprule\noalign{}
\begin{minipage}[b]{\linewidth}\raggedright
Term
\end{minipage} & \begin{minipage}[b]{\linewidth}\centering
Symbol
\end{minipage} & \begin{minipage}[b]{\linewidth}\raggedright
Definition
\end{minipage} & \begin{minipage}[b]{\linewidth}\centering
First Use
\end{minipage} \\
\midrule\noalign{}
\endhead
Cognitively privileged representation & --- & Representation format that
leverages perceptual and spatial cognition for inference & §1 \\
Diagram depth & --- & Length of the longest input-to-output path through
boxes; measures derivational complexity & §4 \\
Existential graphs & --- & Peirce's graphical logic system conducting
first-order logic entirely diagrammatically & §7 \\
Free ride & --- & Shimojima's term: information extracted from a diagram
without explicit inference steps & §1 \\
Hybrid reasoning & --- & Giardino's term: reasoning combining perceptual
pattern recognition with theoretical knowledge & §1 \\
Inferential instrument & --- & Manders's term: a diagram whose spatial
properties encode proof-relevant information & §6 \\
Joyal--Street theorem & --- & Soundness and completeness of
string-diagrammatic reasoning for monoidal categories & §3 \\
Normal form & \(D_{\text{nf}}\) & Canonical form of a diagram obtained
by rewriting; unique up to the axioms & §4 \\
\bottomrule\noalign{}
\end{longtable}
}
\end{frame}

\begin{frame}[fragile]{J. Notation Conventions}
\protect\phantomsection\label{j.-notation-conventions}
{\def\LTcaptype{none} % do not increment counter
\begin{longtable}[]{@{}
  >{\raggedright\arraybackslash}p{(\linewidth - 2\tabcolsep) * \real{0.5000}}
  >{\raggedright\arraybackslash}p{(\linewidth - 2\tabcolsep) * \real{0.5000}}@{}}
\toprule\noalign{}
\begin{minipage}[b]{\linewidth}\raggedright
Convention
\end{minipage} & \begin{minipage}[b]{\linewidth}\raggedright
Meaning
\end{minipage} \\
\midrule\noalign{}
\endhead
\(\mathcal{C}, \mathcal{D}\) & Categories \\
\(\mathcal{V}\) & Enrichment base (monoidal category, typically
\(([0,1], \cdot, 1)\)) \\
\(\mathcal{U}\) & Universal (maximal) case category \\
\(\mathcal{L}\) & Language-specific case category \\
\(\mathcal{E}_{\mathbb{T}}\) & Classifying topos of theory
\(\mathbb{T}\) \\
\(f, g, h\) & Morphisms \\
\(F, G\) & Functors \\
\(\alpha, \beta\) & Natural transformations \\
\(n, s\) & Basic pregroup types: noun, sentence \\
\(n^l, n^r\) & Left and right adjoints of type \(n\) \\
\(n_{\text{NOM}}, n_{\text{ACC}}\) & Case-subscripted noun types (e.g.,
nominative noun, accusative noun) \\
\(N, S\) & Noun space and sentence space under the meaning functor \\
\(N_{\text{NOM}}, N_{\text{ACC}}, \ldots\) & Case-specific vector
subspaces in case-enriched DisCoCat \\
\(\overrightarrow{\text{word}}\) & Word vector (column vector in noun
space \(N\)) \\
\(\overleftrightarrow{\text{verb}}\) & Verb tensor (element of
\(N \otimes S \otimes N\) for transitive verbs) \\
\(\mathbf{Preg}\) & Category of pregroup types and reductions \\
\(\mathbf{FVect}\) / \(\mathbf{FdVect}\) & Category of
finite-dimensional vector spaces and linear maps \\
\(\mathbf{FHilb}\) & Category of finite-dimensional Hilbert spaces and
linear maps \\
\(\mathbf{Qubit}\) & Category of qubit systems with tensor product
structure \\
\(\mathbf{Set}\) & Category of sets and functions \\
\(\otimes\) & Tensor product (monoidal product, type concatenation) \\
\(\circ\) & Composition of morphisms \\
\(\Rightarrow\) & Natural transformation between functors \\
\(\simeq\) & Categorical equivalence \\
\(\leq\) & Preorder relation on types (derivability) \\
\(\gamma\) & Discount factor in distributional RL return computation \\
\(w \in [0,1]\) & Enriched morphism weight (proto-role satisfaction
degree) \\
\(\varepsilon\) & Prediction error (active inference); also compact
closure counit \\
\(\eta\) & Compact closure unit (cap); also learning rate in some
contexts \\
\(\rho\) & Density operator (quantum state) \\
\([@key]\) & Parenthetical citation \\
\texttt{{[}-@key{]}} & Suppress-author citation \\
\texttt{\textbackslash{}autoref\{...\}} & Automatic cross-reference
(section or figure) \\
\bottomrule\noalign{}
\end{longtable}
}
\end{frame}

\end{document}
